\chapter{Giới Thiệu}
\label{chap:intro}

\section{Đặt Vấn Đề và Động Lực}

Âm nhạc là ngôn ngữ phổ quát của cảm xúc con người.
Tuy nhiên, để tạo ra âm nhạc đòi hỏi nhiều năm luyện tập kỹ năng nhạc cụ
và nền tảng kiến thức nhạc lý sâu rộng --- điều mà phần lớn dân số không có.

Câu hỏi nghiên cứu trung tâm của đồ án này là:
\begin{quote}
\textit{Liệu một người không biết gì về nhạc lý có thể tạo ra âm nhạc có cảm xúc
chỉ bằng cách di chuyển cơ thể tự nhiên trước một chiếc webcam thông thường?}
\end{quote}

Đây là hướng tiếp cận khác biệt so với các hệ thống điều khiển âm nhạc truyền thống
(nhạc cụ MIDI, bàn phím điện tử) ở chỗ: người dùng không cần học bất kỳ quy ước
điều khiển nào --- hệ thống tự học và thích nghi với ngôn ngữ cơ thể của họ.

\section{Mục Tiêu Đồ Án}

\begin{enumerate}
  \item Xây dựng hệ thống sinh âm nhạc real-time từ cử chỉ người dùng
        qua webcam thông thường.
  \item Thiết kế và áp dụng kiến trúc hai nhánh:
        \textbf{Gesture Emotion Encoder} (Transformer) và
        \textbf{Music Prior} (GPT), kết hợp qua Cross-Attention.
  \item Tự thiết kế giao thức thu thập và xây dựng dataset cặp
        (chuỗi cử chỉ -- nhãn cảm xúc).
  \item Huấn luyện Music Prior từ dữ liệu MIDI thật để đảm bảo
        chất lượng âm nhạc.
  \item Tích hợp đầu ra vào hệ thống tổng hợp âm thanh MIDI
        real-time qua FluidSynth.
\end{enumerate}

\section{Phạm Vi và Giới Hạn}

Hệ thống tập trung vào \textbf{melody đơn} (single-voice melody generation)
và không giải quyết bài toán soạn nhạc đa bè phức tạp.
Camera đơn sắc (webcam thông thường) được sử dụng mà không cần thiết bị đặc biệt.
Dataset được thu thập từ 2 người trong môi trường có kiểm soát.

\section{Đóng Góp Của Đồ Án}

So với các hệ thống tương tự trong tài liệu, đồ án có các đóng góp:
\begin{enumerate}
  \item \textbf{Tách biệt luồng học cảm xúc và ngữ pháp âm nhạc:}
        Không map trực tiếp cử chỉ sang nốt (như các hệ thống rule-based),
        mà học hai không gian độc lập rồi kết hợp qua Cross-Attention.
  \item \textbf{Dataset tự xây dựng:} 1.853 mẫu với 6 chế độ cảm xúc
        được thiết kế có chủ đích dựa trên mô hình Valence-Arousal.
  \item \textbf{Music Prior từ dữ liệu thật:} Huấn luyện trên
        5.000 bản MIDI từ Lakh MIDI Dataset, đảm bảo nhạc sinh ra
        tuân thủ ngữ pháp âm nhạc thực tế.
\end{enumerate}

\section{Bố Cục Báo Cáo}

Báo cáo được tổ chức như sau:
\begin{itemize}
  \item \textbf{Chương~\ref{chap:related}:} Tổng quan các nghiên cứu liên quan.
  \item \textbf{Chương~\ref{chap:method}:} Kiến trúc hệ thống và phương pháp đề xuất.
  \item \textbf{Chương~\ref{chap:data}:} Thu thập, xử lý và phân tích dữ liệu.
  \item \textbf{Chương~\ref{chap:exp}:} Thiết lập thực nghiệm và kết quả đánh giá.
  \item \textbf{Chương~\ref{chap:conclude}:} Kết luận và hướng phát triển.
\end{itemize}
